\documentclass{jsarticle}           % 日本語用
%\documentclass[fleqn,twocolumn,10pt]{jsarticle}
%   [数式左揃え,ページ分割,フォントサイズ]{日本語用}
%\setlength{\mathindent}{1zw}
%\setlength{\columnsep}{3zw}

%===============================================================================
% section の階層の深さ
%===============================================================================
% LaTeXコマンドシート一覧 - 部・章・節・小節・小々節
% http://www002.upp.so-net.ne.jp/latex/bu.html
% \part*{名前}
% \chapter*{名前}
% \section*{名前}
% \subsection*{名前}
% \subsubsection*{名前}
%
% sectioning - How to add an extra level of sections with headings below \subsubsection - TeX - LaTeX Stack Exchange
% https://tex.stackexchange.com/a/60212
\usepackage{titlesec}
\setcounter{secnumdepth}{3}
\setcounter{tocdepth}{3}
% sectioning - Creating a bold and indented title. - TeX - LaTeX Stack Exchange
% https://tex.stackexchange.com/a/78775
%\renewcommand\thesection{\Roman{section}} 
%\renewcommand\thesubsection{\thesection.\Roman{subsection}} 
\renewcommand\thesection{\arabic{section}} 
\renewcommand\thesubsection{\thesection.\arabic{subsection}} 
\renewcommand\thesubsubsection{\thesubsection.\arabic{subsubsection}} 
% Font Size
% \part*{名前}
% \chapter*{名前}
% \section*{名前}
% \subsection*{名前}
% \subsubsection*{名前}
%\titleformat{\part         }[block]{\Huge }{\thepart.        }{1em}{} 
%\titleformat{\chapter      }[block]{\huge }{\thechapter.     }{1em}{} 
\titleformat{\section      }[block]{\LARGE}{\thesection.     }{1em}{} 
\titleformat{\subsection   }[block]{\Large}{\thesubsection   }{1em}{}
\titleformat{\subsubsection}[block]{\large}{\thesubsubsection}{1em}{}

%===============================================================================
%===============================================================================
\usepackage{listings}               % ソースコード用
\usepackage[dvipdfmx]{graphicx}     % 図の挿入用
%===============================================================================
%===============================================================================
\usepackage{amsmath,amssymb,amsthm} % 数式
% definition
\theoremstyle{plain}
\newtheorem{thm}{Theorem}[section]
\newtheorem{prop}[thm]{Proposition}
\theoremstyle{definition}
\newtheorem{defn}[thm]{Definition}
\theoremstyle{remark}
\newtheorem{rem}{Remark}[thm]
%===============================================================================
%===============================================================================
\usepackage{float}                % 図・表位置[H]
\usepackage{subcaption}           % subtable可
\usepackage[dvipdfmx,]{hyperref}  % Link/URL
\usepackage{pxjahyper}            % "hyperref"の日本語文字化け対策
%\usepackage{url}
\usepackage{slashbox}             % 表の左上の斜線等 (sty file)
    % download to the same directory
    % ref : https://www.biwako.shiga-u.ac.jp/sensei/kumazawa/tex/slashbox.html
\usepackage{color,soul}           % highlight (sty file)
\usepackage{makeidx}              % 索引
\makeindex                        % 索引
\usepackage{tcolorbox}
\usepackage{ascmac}
\usepackage{enumitem}
\usepackage{tikz-network}

%===============================================================================
%===============================================================================

% センターライン(twocolumn)
%\setlength{\columnseprule}{0.4pt}

% 行列matrixのサイズをスケーリングする
\newcommand\scalemath[2]{\scalebox{#1}{\mbox{\ensuremath{\displaystyle #2}}}}

%===============================================================================
%===============================================================================
% ハイライト
%% 日本語ハイライト
\newcommand{\HL}[1]{\hl{\mbox{#1}}}
%% ref: http://otoya8bit.hatenablog.jp/entry/2013/11/28/153226
%% カラー設定
%\sethlcolor{blue}

%===============================================================================
% url custom  
%===============================================================================
% hyperref
\hypersetup{
    colorlinks,  % これをtrueにすると何故かbuildしたとき画像(image)がpdfに表示されなくなったことがあった
    linkcolor=blue,
    filecolor=magenta,      
    urlcolor=cyan,
    bookmarksopen,
    bookmarksdepth=2,
    bookmarkstype=toc,  % toc link
%    linktocpage=true,   % tocのページ番号にリンクが現れる
    bookmarks=true,     % toc link
    % ref
    % https://texwiki.texjp.org/?hyperref#h026c44c
}


%===============================================================================
% tableofcontents
%===============================================================================
\makeatletter
    %\renewcommand{\l@<typ>        }{\@dottedtocline{<level>}{<indentation>}{<numwidth>}}
    \renewcommand*{\l@subsection   }{\@dottedtocline{      2}{         10em}{       5em}}
    \renewcommand*{\l@subsubsection}{\@dottedtocline{      3}{         15em}{       5em}}
\makeatother
        %The command \@dottedtocline expects the following parameters:
        %\renewcommand{\l@<typ>}{\@dottedtocline{<level>}%
        %                                       {<indentation>}%
        %                                       {<numwidth>}}
        % https://tex.stackexchange.com/questions/33841/how-to-modify-the-space-between-the-numbers-and-text-of-sectioning-titles-in-the/33842


%===============================================================================
%===============================================================================
\title{タイトル}
\author{pollenJP}

%===============================================================================
%===============================================================================
\begin{document}
    
\maketitle
    
\renewcommand\contentsname{目次（もくじだよ）}
\tableofcontents
\addtocontents{toc}{
  {\fontsize{20.0pt}{0pt}\selectfont 目次}
  \hfill\textbf{Page}\par}
\addcontentsline{toc}{section}{Preface}
\addcontentsline{toc}{subsection}{Preface}
    
\section{TableOfContents}
\begin{itemize}
  \item \url{https://tex.stackexchange.com/questions
              /33841/how-to-modify-the-space-between-the-numbers-
              and-text-of-sectioning-titles-in-the/33842
        }
  \item \url{https://en.wikibooks.org/wiki/LaTeX/Document_Structure}
  \item \url{https://cmry.github.io/notes/latextoc}
\end{itemize}

\section{環の定義}
\begin{enumerate}[label=\fbox{\arabic*}]
  \item
        {[問]}\\
        $R$を環とする. $0 \in R$を零元とすると, $\forall a \in R$ に対して,
        \[ a \cdot 0 = 0 \cdot a = 0 \] となることを示せ. \\
        {[解答]}\\
        \begin{align*}
                          & 0 = 0 + 0                                                 \\
                          & (\because 零元の定義)                                \\
          \Leftrightarrow & a \cdot 0 = a \cdot ( 0 + 0 )                             \\
                          & (\text{左からaをかける})                           \\
          \Leftrightarrow & a \cdot 0 = a \cdot 0 + a \cdot 0                         \\
                          & (\because \text{分配則})                               \\
          \Leftrightarrow & a \cdot 0 - a \cdot 0 = a \cdot 0 + a \cdot 0 - a \cdot 0 \\
                          & (\because \text{両辺に加法逆元を加える})       \\
          \Leftrightarrow & 0 = a \cdot 0                                             \\
                          & (\because \text{結合則と加法逆元の性質})       
        \end{align*}
        同様に
        \begin{align*}
          \Leftrightarrow & 0 \cdot a = (0 + 0) \cdot a                               \\
                          & (\text{右からaをかける})                           \\
          \Leftrightarrow & 0 \cdot a = 0 \cdot a + 0 \cdot a                         \\
                          & (\because \text{分配則})                               \\
          \Leftrightarrow & 0 \cdot a - 0 \cdot a = 0 \cdot a + 0 \cdot a - 0 \cdot a \\
                          & (\because \text{両辺に加法逆元を加える})       \\
          \Leftrightarrow & 0 = 0 \cdot a                                             \\
                          & (\because \text{結合則と加法逆元の性質})       
        \end{align*}
        よって, $a \cdot 0 = 0 \cdot a = 0$ \hfill $\Box$
        
  \item
        {[問]}\\
        環$R$において乗法の単位元は一意であることを示せ.\\
        {[解答]}\\
        $e^{*}, e^{'*} \in R$がともに環$R$の乗法の単位元であるとする.
        \begin{align*}
          e^{*} & = e^{*} \cdot e^{'*} & (\because e^{'*} \text{は単位元}) \\
                & = e^{'*}             & (\because e^{*} \text{は単位元})  
        \end{align*}
        \hfill $\Box$
        
  \item
        {[問]}\\
        $R$を環とする. $a \in R$に対して左逆元$b$と右逆元$c$が存在するとき, $b=c$であることを示せ.\\
        また,それらを$a$の逆元というが, $a$の逆元が一意であることを示せ.\\
        {[解答]}\\
        \begin{itembox}[l]{左逆元・右逆元の定義}
          hello
        \end{itembox}
\end{enumerate}

    
\section{Who is pollenJP?}
\begin{tcolorbox}
  hhh
\end{tcolorbox}

\subsection{list}
\subsubsection{pollenJPJPJP}
\begin{enumerate}
  \item \HL{ハイライト}
        \sethlcolor{red}
  \item \HL{カラー変更}
  \item ccc\index{ccc}
  \item ddd
\end{enumerate}
\begin{enumerate}
  \setcounter{enumi}{8}
  \item aaa
  \item bbb
  \item ccc
  \item \url{https://tex.stackexchange.com/questions
              /142/how-can-i-make-an-enumerate-list-start
        -at-something-other-than-1}
\end{enumerate}
%
\begin{description}
  \item[Block のパラメータの設定 Gaussian Noise Generator]\mbox{}
  \begin{itemize}
    \item 平均値(Mean Value)=0
    \item 分散(Variance)=0.8
    \item サンプル時間=0.01 秒
  \end{itemize}
  \item[シミュレーション時間]\mbox{}
  \begin{itemize}
    \item 終了時間 10 秒
  \end{itemize}
  \item[その他]\mbox{}
  \begin{itemize}
    \item Signal To Workspace,Display は,設定の変更は必要ない.
  \end{itemize}
  \item[hello]\mbox{}\\
  あいうえお
\end{description}

\subsubsection{pollenJPJPJP2}
hello
\subsection{pollenJPJP2}
hello
    
\section{表}
\subsection{Normal}
以下の表\ref{課題1.2: M=2の表}を示す.
% table
\begin{table}[H]  % t(top),b(bottom),p(page),h(できればhere),H(必ずhere)
  \caption[M=2]{課題1.2: M=2の表}
  \label{課題1.2: M=2の表}
  \begin{center}
    \setlength{\tabcolsep}{3pt}
    \footnotesize
    \begin{tabular}{|c|c|c} \hline
      シンボル & ビット \\ \hline \hline
      0            & 0         \\ \hline
      1            & 1         \\ \hline
    \end{tabular}
  \end{center}
\end{table}
        
\subsection{subtable}
以下の表\ref{課題1.2: 課題}
% table
\begin{table}[H]
  \begin{subtable}{0.49\columnwidth}
    \begin{center}
      \caption[M=2]{M=2の表}
      \label{課題1.2: M=2の表}
      \setlength{\tabcolsep}{3pt}
      \footnotesize
      \begin{tabular}{|c|c|c}
        \hline
        シンボル & ビット \\ \hline \hline
        0            & 0         \\ \hline
        1            & 1         \\ \hline
      \end{tabular}
    \end{center}
  \end{subtable}
  \begin{subtable}{0.49\columnwidth}
    \begin{center}
      \caption[M=4]{M=4の表}
      \label{課題1.2: M=4の表}
      \setlength{\tabcolsep}{3pt}
      \footnotesize
      \begin{tabular}{|c|c|c}
        \hline
        シンボル & ビッ�� \\ \hline \hline
        0            & 00        \\ \hline
        1            & 01        \\ \hline
        2            & 10        \\ \hline
        3            & 11        \\ \hline
      \end{tabular}
    \end{center}
  \end{subtable}
  \begin{subtable}{\columnwidth}
    \begin{center}
      \caption[M]{課題}
      \label{課題1.2: 課題}
      \setlength{\tabcolsep}{3pt}
      \footnotesize
      \begin{tabular}{|l|c|c|c|c|c|c|c|c|c|c|c|c|}
        \hline
                                    
        & \multicolumn{6}{|l|}{M=2}
        & \multicolumn{6}{|l|}{M=4} \\ \hline
        時間 
          & 0  & 1  & 2  & 3  & 4  & 5  
          & 0  & 1  & 2  & 3  & 4  & 5  \\ \hline
        シンボル列
          & 1  & 0  & 0  & 0  & 1  & 1  
          & 2  & 0  & 0  & 1  & 2  & 3  \\ \hline
        ビット列
          & 1  & 0  & 0  & 0  & 1  & 1  
          & 10 & 00 & 00 & 01 & 10 & 11 \\ \hline
      \end{tabular}
    \end{center}
  \end{subtable}
\end{table}
            
\subsection{斜線}
% table
\begin{table}[H]
  \begin{center}
    \caption{課題2.1: 雑音の統計値(平均値と分散)}
    \label{課題2.1: 雑音の統計値(平均値と分散)}
    \begin{tabular}{|l||c|c|}
      \hline
      \backslashbox{統計値}{分散(設定値)} & 1       & 2       \\ \hline\hline
      平均値                                   & -0.0247 & -0.0349 \\ \hline
      分散                                      & 0.9850  & 1.97    \\ \hline
    \end{tabular}
  \end{center}
\end{table}

\subsection{$p$次元データ（共分散・相関係数）}
\HL{\textgt{偏差積和}}
\begin{align*}
  S_{ij} = \frac{1}{n} \sum\limits_{i=1}^n           
  (x_{j_1j_2}-\bar{x_{j_1}})(x_{ij_2}-\bar{x_{j_2}}) 
\end{align*}
% table
\begin{table}[H]
  \raggedright
  \caption{多変量データ例}
  \label{第2回 多変量データ例}
  \scalebox{0.8}[1.0]{
    \begin{tabular}{|r|c|c|c|} \hline
      No.      & $x_{j_1}$ & $x_{j_2}$ & 偏差積 \\
      \hline \hline
      $1$ & $x_{1j_1}$ & $x_{1j_2}$ 
      & $(x_{1j_1}-\bar{x_{j_1}})(x_{1j_2}-\bar{x_{j_2}})$
      \\ \hline
      $2$ & $x_{2j_1}$ & $x_{2j_2}$ 
      & $(x_{2j_1}-\bar{x_{j_1}})(x_{2j_2}-\bar{x_{j_2}})$
      \\ \hline
      $\vdots$ &           &           &           \\ \hline
      $i$ & $x_{ij_1}$ & $x_{1j_2}$ 
      & $(x_{ij_1}-\bar{x_{j_1}})(x_{ij_2}-\bar{x_{j_2}})$
      \\ \hline
      $\vdots$ &           &           &           \\ \hline
      $n$ & $x_{nj_1}$ & $x_{nj_2}$ 
      & $(x_{nj_1}-\bar{x_{j_1}})(x_{nj_2}-\bar{x_{j_2}})$
      \\ \hline
      & $\bar{x_{j_1}}$ &  $\bar{x_{j_2}}$ 
      & $S_{ij} = \frac{1}{n} \sum\limits_{i=1}^n 
      (x_{j_1j_2}-\bar{x_{j_1}})(x_{ij_2}-\bar{x_{j_2}})$
      \\ \hline
    \end{tabular}
  }
\end{table}


    
\section{図} \label{section図}
以下の図
%figure
\begin{figure}
          
\end{figure}
% figure
%    \begin{figure}[H]  % t(top),b(bottom),p(page),h(できればhere),H(必ずhere)
%      \begin{center}
%        \includegraphics[width=0.8\textwidth]{fig01_kadai1-1.png}
%        \caption{課題1.1:構成}
%        \label{課題1.1:構成}
%      \end{center}
%    \end{figure}
%

%%%%%%%%%%%%%%%%%%%%%%%%%%%%%%%%%%%%%%%%
\section{TIKZ-NETWORK - 描画} \label{TIKZ}
\begin{itemize}
  \item ここにいろいろ公開されている \url{https://arxiv.org/abs/1709.06005}
  \item tikz-network.styをダウンロードしなければいけない
        \begin{itemize}
          \item \url{https://ctan.org/tex-archive/graphics/pgf/contrib/tikz-network}
        \end{itemize}
\end{itemize}
\begin{tikzpicture}
  \filldraw ( -.2, .2) circle (2pt) (.2,.2) circle (2pt);
  \draw (0,0) circle (5mm) (-.3, -.1) .. controls (0,-.3) .. (.3,-.1);
\end{tikzpicture}
%%%%%%%%%%%%%%%%%%%%%%%%%%%%%%%%%%%%%%%%

    
% save old section configure
\let\oldthesection=\thesection
\let\oldthesubsection=\thesubsection
%  \setcounter{section}{0}
\renewcommand{\thesection}{課題\arabic{section}}
\renewcommand{\thesubsection}{\thesection.\arabic{subsection}}
\renewcommand{\thesubsubsection}{\thesubsection.\arabic{subsubsection}}
%\renewcommand{\thesubsubsection}{課題\arabic{section}.\arabic{subsection}}
\section{セクションタイトル替え} \label{課題1}
\subsection*{実験内容}
\subsection{サブタイトルも変える}
\subsubsection{subsub}
\begin{itemize}
  \item \url{
        https://groups.google.com/forum/#!topic/comp.text.tex/EcAPeYr-ySE}
\end{itemize}
\subsubsection*{subsub2}
\subsection{次のセクションでタイトルを戻す}

\let\thesection=\oldthesection
\let\thesubsection=\oldthesubsection
\section{URLを使用}
\begin{itemize}
  \item \url{https://www.sharelatex.com/learn/Hyperlinks}
\end{itemize}
    
\section{数式} \label{数式}
\subsection{Show inline math as if it were display math}
\url{https://tex.stackexchange.com/questions/32824/
show-inline-math-as-if-it-were-display-math}
$\frac{1}{n-1} \sum\limits_{i=1}^{n} (x_{ij} - \bar{x_j})^2$
            
$\prod\limits_{j=0}^J k_j$
            
$\int\limits_{-\infty}^\infty f(x)\,\mathrm{d}x$
            
\subsection{実験内容}
変調とは,伝送媒体である波(電波,光など)に情報を乗せる操作である.
波は.振幅と周波数,位相をパラメータとして構成される.
伝送媒体の波を搬送波(Carrier)と呼ぶ.搬送波は,次式で表せる.
% math
\begin{align}
  \label{math:搬送波}                              
  s(t) = r(t) \cdot cos(2 \pi f_c t + \vartheta (t) ) 
\end{align}
$f_c$ を搬送波周波数と呼ぶ.携帯電話の代表値では 860MHz である.
式(\ref{math:搬送波})を展開すると次式を得る.
% math
\begin{align}
  \label{math:搬送波（展開）}                               
  s(t) = I(t) \cdot cos(2 \pi f_c t) - Q(t) \cdot sin(2 \pi f_c t) 
\end{align}
式(\ref{math:搬送波})は式(\ref{math:搬送波（展開）})と等価であり,
表現の方法が違うのみである.式(\ref{math:搬送波（展開）})の表現において
$I(t)$を搬送波の同相成分,$Q(t)$を直交成分と言う.
これは余弦波と正弦波が互いに直交(1 周期積分してゼロとなる)しているためである.
さらに,式(\ref{math:搬送波（展開）})は次のように表せる.
% math
\begin{align}
  \label{搬送波（展開）別表現}
  \begin{split}
  s(t) 
       & = Re \left\{ \left( I(t) + j \cdot Q(t) \right)   
  \cdot e^{j2 \pi f_c t} \right\} \\
       & = Re \left\{ u(t) \cdot e^{j2 \pi f_c t} \right\} \\
  u(t) & = I(t) + j \cdot Q(t)                             
  \end{split}
\end{align}

\subsection{行列のスケーリング}
% 行列matrixのサイズをスケーリングする
%\newcommand\scalemath[2]{\scalebox{#1}{\mbox{\ensuremath{\displaystyle #2}}}}
\begin{align*}
  X
  &= \begin{bmatrix}
  \sum\limits_{i=1}^n (x_{i1}-\bar{x_1})^2 & \cdots 
  & \sum\limits_{i=1}^n (x_{i1}-\bar{x_1})(x_{ip}-\bar{x_p}) \\
  \vdots & \ddots & \vdots                                   \\
  \sum\limits_{i=1}^n (x_{ip}-\bar{x_p})(x_{i1}-\bar{x_1}) 
         & \cdots & \sum\limits_{i=1}^n (x_{ip}-\bar{x_p})^2 
  \end{bmatrix} \\
  &= \scalemath{0.3}{
  \begin{bmatrix}
  \sum\limits_{i=1}^n (x_{i1}-\bar{x_1})^2 & \cdots 
  & \sum\limits_{i=1}^n (x_{i1}-\bar{x_1})(x_{ip}-\bar{x_p}) \\
  \vdots & \ddots & \vdots                                   \\
  \sum\limits_{i=1}^n (x_{ip}-\bar{x_p})(x_{i1}-\bar{x_1}) 
         & \cdots & \sum\limits_{i=1}^n (x_{ip}-\bar{x_p})^2 
  \end{bmatrix}
  }
\end{align*}

    
\section{プログラミングコード}
以下のPython3コードによって描画した.
% Coding
\begin{lstlisting}[basicstyle=\ttfamily\footnotesize, frame=single]
import numpy as np
import matplotlib.pyplot as plt

f = np.arange(-5.0, 5.1, 0.1)
print(f.shape)

T = 2
G_f = T * np.sin(np.pi * f * T)
W_f = abs(G_f)**2 / T
P_f = 10 * np.log10(W_f)

fig = plt.figure(figsize=(15, 10))
ax = fig.add_subplot(1,1,1)
ax.plot(f,P_f, label="P(f)")
ax.set_title("Power Spectrum")
ax.set_xlabel("frequency[Hz]")
ax.set_ylabel("Magnitude-squared [dB]")
ax.legend(loc="lower right", prop={'size': 20})
plt.show()
\end{lstlisting}
        


        

\printindex
\end{document}

        
        

